(1) 将入射速度分解为 $\mathbf v=\mathbf v_t+v_n\hat n$，其中法向单位向量指向来球一侧。恢复系数定义使碰后法向分量 $v_n'=-ev_n$，光滑壁无切向冲量，因此 $\mathbf v_t'=\mathbf v_t$，总速度为 $\mathbf v'=\mathbf v_t-ev_n\hat n$。
(2) 冲量由运动量改变量得到 $\mathbf J=m(\mathbf v'-\mathbf v)=-m(1+e)v_n\hat n$。由于 $v_n<0$，这个矢量指向壁外、远离壁面，模为 $J=m(1+e)|v_n|$，单位为千克米每秒。
(3) 切向动能不变，法向动能由 $mv_n^2/2$ 变为 $me^2v_n^2/2$。因此碰撞损失能量 $\Delta K=m(1-e^2)v_n^2/2\ge0$；当 $e<1$ 且碰撞具有非零法向速度时该能量损失严格大于零。
(4) $e=0$ 时仅法向分速度消失，切向滑行依然存在；$e=1$ 是法向镜面弹性反射；$v_n\to0$ 时冲量与损耗都趋零。若接触面粗糙或球存在相关的自旋摩擦，则切向冲量未必为零，需要进一步加入角动量方程。
